\documentclass[12pt]{article}
\usepackage[utf8]{inputenc}
\usepackage{amsmath}
\usepackage{amssymb}
\usepackage{geometry}
\usepackage{hyperref}
\usepackage{booktabs} % for nicer tables
\usepackage{float}
\usepackage{graphicx}
\usepackage{longtable}
\usepackage{pdfpages}
\usepackage{url}
\geometry{a4paper, margin=1in}

\title{Nichols Radial Injection Model (RIM) VI The Birth of the Universe}
\author{Lawrence William Nichols}
\date{\today}

\begin{document}

\maketitle

\begin{abstract}
A Description of the Formation of Stars, Planets, and Black Holes, plus there alignments, positions and orbits.
\end{abstract}

\section{Introduction}
In the beginning there was the first injection of matter from the hypersphere via gamma ray bursts, this contain enough material to create dust clouds with hydrogen and heavier elements. The first
particles joined together and as these grew larger they eventually collected enough matter to created stars and black holes these became the first primitive solar systems. 

\section{Black Hole Accretion and Magnetic Ejection}

Black holes gather matter through an accretion disk. The accretion rate can be limited by the Eddington luminosity:
\begin{equation}
    L_{\rm Edd} = \frac{4\pi G M m_p c}{\sigma_T} \approx 1.26 \times 10^{38} \left( \frac{M}{M_\odot} \right) \, \rm erg\,s^{-1},
\end{equation}
where $M$ is the black hole mass, $m_p$ the proton mass, and $\sigma_T$ the Thomson cross-section. This sets an approximate upper limit on the mass accretion rate $\dot{M} \approx L_{\rm Edd}/(\eta c^2)$ with efficiency $\eta \sim 0.1$.

Magnetic vortices above the disk can eject material via magnetohydrodynamics processes. The magnetic torque on accreting material scales roughly as
\begin{equation}
    \tau_B \propto B^2 r^3,
\end{equation}
where $B$ is the magnetic field strength and $r$ the radial distance. Iron, being highly magnetic (high susceptibility $\chi$), is more easily captured and accelerated along field lines. These ejected masses may clump together and eventually form planetary bodies when captured by stars.

\section{Planetary Orbit Relationship}

An interesting observation from studying our solar system is that planetary orbits may correlate with the size of the iron core of the planet. For example, Mercury has the largest relative iron core (core mass fraction CMF $\approx 0.7$ by mass, or $\sim 85\%$ by volume) and the closest orbit, while Pluto has one of the smallest and the most distant.

In the Nichols Radial Injection Model, we hypothesize an inverse relationship between the iron core mass fraction $f_{\rm Fe}$ and semi-major axis $a$:
\begin{equation}
    f_{\rm Fe}(a) \propto \frac{1}{a^\alpha} \quad \text{or} \quad f_{\rm Fe}(a) \approx k \exp\left(-\frac{a}{a_0}\right),
\end{equation}
where $\alpha > 0$ and $a_0$ is a characteristic length scale related to magnetic sorting or injection dynamics in the early protoplanetary environment. For the solar system, Mercury ($a \approx 0.39$ AU, high $f_{\rm Fe}$) and outer bodies (low $f_{\rm Fe}$) suggest $\alpha \sim 1$--$2$ as a first approximation (though this is highly schematic and requires further testing).

\section{Nichols Radial Injection Model}

In the Nichols Radial Injection Model (RIM), the first injection of matter came from a gamma-ray burst similar in scale to GRB~250702~BDE. This event injected enough mass $M_{\rm inj}$ to form the first black hole, starting the cycle of star and planetary formation.

The injected mass must exceed the threshold for direct collapse to a black hole:
\begin{equation}
    M_{\rm inj} > M_{\rm crit} \approx \frac{c^2 R_S}{2G} = M_{\rm BH},
\end{equation}
where $R_S = 2GM/c^2$ is the Schwarzschild radius of the resulting black hole. Subsequent accretion and ejection cycles follow:
\begin{equation}
    \frac{dM_{\rm BH}}{dt} = \dot{M}_{\rm acc} - \dot{M}_{\rm ej},
\end{equation}
with ejected material $\dot{M}_{\rm ej}$ feeding star and planet formation in orbiting disks.

\section{Black Hole Particle Accelerator}
\url{https://www.youtube.com/watch?v=rh6XaL_W1ns&t=99s}

\section{Conclusion}
Over billions of years this cycle of black hole formation, stellar birth,
and planetary creation could produce the universe as observed today.

\begin{thebibliography}{99}

\bibitem{RIM-I}
L.\,W.\,Nichols.
Nichols Radial Injection Model (RIM) and Radial ERB Inflows: A Mechanism for Progenitor-less Astrophysical Events, the Hubble Pulse, and 4D Substrate Interactions.
[Self-published], January 21, 2026.
\newline
\url{https://rxiverse.org/abs/2602.0036}

\bibitem{RIM-II}
L.\,W.\,Nichols.
Nichols Radial Injection Model (RIM) 11-Year Solar Cycle Pulses, GRB Milestones, Historical Alignments, and the 16.6 Gyr Hypersphere Universe.
[Self-published], February 2026.
\newline
\url{https://rxiverse.org/abs/2602.0038}

\bibitem{RIM-III}
L.\,W.\,Nichols.
Nichols Radial Injection Model (RIM) III: The 3,000-Year Universal Odometer, the 1054 CE Primary Strike, and the Decaying 11-Year Cycle.
[Self-published], 15 February 2026.
\newline
\url{https://rxiverse.org/abs/2602.0043}

\bibitem{RIM-IV}
L.\,W.\,Nichols.
Nichols Radial Injection Model (RIM) IV: A 5D Mass-Regulated Manifold Solution to Cosmic Expansion.
[Self-published], 28 February 2026.
\newline
\url{https://rxiverse.org/abs/2603.0001}

\bibitem{RIM-V}
L.\,W.\,Nichols.
Nichols Radial Injection Model (RIM) V:
The Empirical Seal of the 11.07-Year Pulse
and Universal Resonance in Stellar Cycles
[Self-published], 4 March 2026.
\newline
\url{https://rxiverse.org/abs/2603.0003}

\bibitem{gompertz2025}
Gompertz, B.\ P., et al.,
``JWST Spectroscopy of GRB 250702B,''
arXiv:2509.22778, 2025.
\newline
\url{https://arxiv.org/abs/2509.22778}

\bibitem{sn1054_wiki}
Wikipedia,
``SN 1054,''
2026.
\newline
\url{https://en.wikipedia.org/wiki/SN_1054}

\bibitem{egeland2015} Egeland, R., et al. (2015).
``Long-term magnetic activity in the chromosphere of the Sun-like star HD 30495.''
The Astrophysical Journal, 811(2), 111.
\url{https://doi.org/10.1088/0004-637X/811/2/111}

\bibitem{einstein1935}
Einstein, A., \& Rosen, N. (1935).
\textit{The Particle Problem in the General Theory of Relativity}.
Physical Review, 48(1), 73.

\end{thebibliography}

\newpage

\section{Engineering Studies and Visual Data}
This section includes relevant visual diagrams.
\includepdf[page=-]{Core Size.png}
\includepdf[pages=-]{Cosmic_Formation_Blueprint.PDF}

\end{document}
